%% Exponentials, sieve classifier, and power objects for a small presheaf base.

\subsection{Exponentials of presheaves}
\label{sec:p10-presheaf-exponentials}

Fix the small category \(\mathcal C\) and fixed-universe convention of
Definition~\ref{def:p10-presheaf-category}.
Write \(\widehat{\mathcal C}=[\mathcal C^{\mathrm{op}},\mathsf{Set}]\)
and \(\mathrm y c=\mathcal C(-,c)\).
Products are the pointwise products of
Theorem~\ref{thm:p10-presheaf-limits}.
The presheaves \(0\) and \(1\) have respectively the empty set
and a fixed singleton as every component.

\begin{hbtheorem}{thm:p10-presheaf-exponentials}{proved=annex:p10.psh.3}{The exponential presheaf and its evaluation}
For presheaves \(X,Y\), define
\[
                    (Y^X)(c)=
                    \operatorname{Nat}(\mathrm y c\times X,Y).
\]
For \(u:d\to c\), its restriction is precomposition:
\[
 (Y^X)(u)(\alpha)=\alpha\circ(\mathrm y u\times1_X),\qquad
 [(Y^X)(u)(\alpha)]_e(v,x)=\alpha_e(u\circ v,x)
\]
for \(v:e\to d\) and \(x\in X(e)\).
This defines a presheaf with small component sets.
Its evaluation \(\operatorname{ev}:Y^X\times X\to Y\) has components
\[
                         \operatorname{ev}_c(\alpha,x)
                                      =\alpha_c(1_c,x).
\]
For \(h:Z\times X\to Y\), its unique transpose
\(\Lambda h:Z\to Y^X\) is given by
\[
             [(\Lambda h)_c(z)]_d(u,x)=h_d(Z(u)(z),x)
               \quad(u:d\to c,\ z\in Z(c),\ x\in X(d)).
\]
These are natural transformations.  For every \(k:Z\to Y^X\),
the inverse identities are
\[
 \operatorname{ev}\circ(\Lambda h\times1_X)=h,\qquad
 \Lambda\bigl(\operatorname{ev}\circ(k\times1_X)\bigr)=k.
\]
Thus evaluation and transposition give inverse bijections
\[
 \widehat{\mathcal C}(Z\times X,Y)
                   \cong\widehat{\mathcal C}(Z,Y^X),
\]
natural in \(Z\) and \(Y\).  In particular
\(\widehat{\mathcal C}\) is cartesian closed.
The construction includes empty component sets and the empty base.
\end{hbtheorem}

This is the exponential construction of
\cite[I\S6, Proposition~1, pp.~46--47, formulas (5)--(8)]{MacLaneMoerdijk1992}.
The same component and restriction formulas occur in
\cite[\S1, pp.~4--5]{MoerdijkVanOosten2007Topos}.
The element \(\alpha\in(Y^X)(c)\) specifies compatible functions
at every object mapping into \(c\); its data are not in general
just a function \(X(c)\to Y(c)\).

\subsection{The sieve classifier}
\label{sec:p10-presheaf-classifier}

\begin{hbdefinition}{def:p10-presheaf-sieves}{Sieves and the presheaf of truth values}
A sieve \(S\) on \(c\) is a subset of the arrows with codomain \(c\)
closed under precomposition: if \(u:d\to c\) belongs to \(S\)
and \(v:e\to d\), then \(u\circ v\in S\).
Let \(\Omega(c)\) be the set of all sieves on \(c\).
For \(u:d\to c\), put
\[
             \Omega(u)(S)=
                \{v:e\to d:u\circ v\in S\}.
\]
Write \(t(c)\) for the maximal sieve of all arrows into \(c\).
It is the unique sieve containing \(1_c\).
Define the truth components
\[
                    \mathsf{true}_c:1(c)\to\Omega(c),
                    \qquad \mathsf{true}_c(*)=t(c).
\]
The sieve description is equivalently a subpresheaf of
\(\mathrm y c\): its component at \(d\) is
\(S\cap\mathcal C(d,c)\).
\index{sieve}\index{subobject classifier!presheaf}
\end{hbdefinition}

\begin{hbtheorem}{thm:p10-presheaf-topos}{proved=annex:p10.psh.4}{The sieve classifier makes a presheaf category a topos}
The displayed restrictions make \(\Omega\) a presheaf, and
\(\mathsf{true}:1\to\Omega\) a monomorphism.
For a subpresheaf \(A\subseteq X\), define
\[
       (\chi_A)_c(x)=
              \{u:d\to c:X(u)(x)\in A(d)\}.
\]
This is a natural transformation \(X\to\Omega\), and the square
\[
\begin{array}{ccc}
 A&\longrightarrow&1\\
 \downarrow&&\downarrow{\scriptstyle\mathsf{true}}\\
 X&\xrightarrow{\ \chi_A\ }&\Omega
\end{array}
\]
is a pullback.  It is the unique characteristic map for this
subobject.
For an arbitrary monomorphism \(m:B\to X\), use its canonical
image subpresheaf \(A(c)=m_c[B(c)]\); the same characteristic
map gives a pullback square with \(m\) as its left side.

Hence \(\mathsf{true}\) is a subobject classifier.
Together with the finite limits and exponentials already proved,
this makes \(\widehat{\mathcal C}\) an elementary topos in the
sense of Definition~\ref{def:p10-topos}.
For \(\mathcal C=\varnothing\), it is the degenerate category
with one object and one arrow, and still satisfies this definition.
\end{hbtheorem}

Sieves, their restriction, and the characteristic map are
\cite[I\S4, pp.~37--38, formulas (2)--(8)]{MacLaneMoerdijk1992}.
The characteristic-map uniqueness argument is also given in
\cite[\S1, pp.~6--7]{MoerdijkVanOosten2007Topos}.
Smallness of the base is used throughout: each \(\Omega(c)\)
is a subset of the powerset of the set of arrows into \(c\).

\subsection{Power objects and membership}
\label{sec:p10-presheaf-power-objects}

\begin{hbdefinition}{def:p10-presheaf-power-object}{The power object and its membership relation}
For a presheaf \(W\), write \(\operatorname{Sub}(W)\) for its
subobjects, identified with actual subpresheaves using the
canonical images of
Theorem~\ref{thm:p10-presheaf-mono-epi}.
For \(X\), put
\[
                              P(X)=\Omega^X .
\]
Let \(\in_X\hookrightarrow P(X)\times X\) be the pullback
of \(\mathsf{true}:1\to\Omega\) along evaluation.
It is called the membership subobject.
A power object represents subobjects of products with \(X\):
the following theorem gives that universal property explicitly.
\index{power object!presheaf}\index{membership!power object}
\end{hbdefinition}

\begin{hbtheorem}{thm:p10-presheaf-power-objects}{proved=annex:p10.psh.5}{Classification by a power object}
There are bijections, natural in \(Z\),
\[
 \widehat{\mathcal C}(Z,P(X))
       \cong\widehat{\mathcal C}(Z\times X,\Omega)
       \cong\operatorname{Sub}(Z\times X).
\]
The first is uncurrying, and the second pulls back
\(\mathsf{true}\).  Thus \(k:Z\to P(X)\) classifies exactly
\[
                              (k\times1_X)^*(\in_X).
\]
Conversely, the unique classifying map of a subobject
\(S\hookrightarrow Z\times X\) is \(k=\Lambda(\chi_S)\).
Here a superscript \({}^*\) on subobjects denotes pullback.

At each base object \(c\), classification gives a canonical bijection
\[
                       P(X)(c)\cong
                       \operatorname{Sub}(\mathrm y c\times X).
\]
Under this identification, the restriction along \(u:d\to c\)
sends \(S\subseteq\mathrm y c\times X\) to
\[
 S|u=(\mathrm y u\times1_X)^*S,\qquad
 (S|u)(e)=\{(v,x)\in\mathcal C(e,d)\times X(e):
                                      (u\circ v,x)\in S(e)\}.
\]
The membership components are
\[
 (\in_X)(c)=
   \{(S,x):S\in\operatorname{Sub}(\mathrm y c\times X),\
                         x\in X(c),\ (1_c,x)\in S(c)\},
\]
where the first coordinate uses the same canonical identification.
For \(S\subseteq Z\times X\), its classifying map sends
\(z\in Z(c)\) to the subpresheaf
\[
        S_z(d)=\{(u,x)\in\mathcal C(d,c)\times X(d):
                                      (Z(u)(z),x)\in S(d)\}.
\]
These formulas include \(X=0\) and \(\mathcal C=\varnothing\).
All displayed component and subobject collections are small sets.
\end{hbtheorem}

Power objects as \(\Omega^X\), their natural classification, and
the membership formula are described in
\cite[\S1, pp.~7--8, Exercise~4]{MoerdijkVanOosten2007Topos}.
The annex supplies the restriction, membership, and classification
proofs.

\begin{hbexample}{def:p10-presheaf-not-pointwise}{Exponential and power objects depend on restrictions}
Let \(\mathcal C\) have objects \(0,1\) and one nonidentity arrow
\(a:0\to1\).
Let \(X(0)=\{*\}\), \(X(1)=\varnothing\), with its unique restriction.
Then \(\mathrm y1\) is the terminal presheaf.
There is no natural map \(\mathrm y1\times X\to0\), because its
component at \(0\) would map a singleton to the empty set.
Consequently
\[
                     (0^X)(1)=\varnothing,
\]
whereas the external function set
\(\mathsf{Set}(X(1),\varnothing)\) is a singleton.
Also \(\mathrm y1\times X\) has exactly two subpresheaves,
so \(P(X)(1)\) has two elements, whereas
\(\mathcal P(X(1))=\mathcal P(\varnothing)\) has one.
\end{hbexample}
